\section{Related Work}

\subsection{Code-Generation Benchmarks and LLM Priors}
HumanEval~\cite{chen2021codex} and MBPP~\cite{austin2021program} established that LLMs trained on large code corpora develop strong token-level priors: Python keywords appear in statistically predictable positions. More recent code LLMs (DeepSeek-Coder~\cite{guo2024deepseek}, WizardCoder~\cite{luo2023wizardcoder}) amplify these priors through specialized pretraining, while encoder-decoder models (CodeBERT~\cite{feng2020codebert}, CodeT5~\cite{wang2021codet5}) reinforce them via masked language modeling objectives.

\subsection{Instruction-Following in Code Generation}
Recent benchmarks---CodeIF~\cite{codeif2025} and MultiCodeIF~\cite{multicode2025}---reveal that LLMs underperform on code tasks that require strict adherence to developer-specified constraints, even when those constraints are explicitly stated. These works show a gap between general programming competence and instruction compliance, motivating our study of a more extreme case: instructions that directly conflict with Python NTP priors.

\subsection{Knowledge Conflicts: Parametric vs.\ Contextual}
A growing body of work studies how LLMs resolve conflicts between pretrained parametric knowledge and information given in context.
Chen et al.~\cite{chen2024instructhierarchy} propose an instruction hierarchy to prioritize privileged instructions; however, they focus on natural-language directives rather than programming-language semantic rules.
Xie et al.~\cite{xie2024conflict} and Wu et al.~\cite{wu2024contextparametric} show that LLMs often default to parametric knowledge when context contradicts it---especially under low surface-form overlap---which directly motivates our L3/L4 designs where syntax is identical to Python but semantics are inverted.
Wang et al.~\cite{wang2024ntp} characterize how NTP training shapes model behavior; their ``law of next-token prediction'' provides theoretical grounding for our observation that prior dominance scales with token predictability.

\subsection{Adversarial Perturbations in Code Generation}
PromptBench~\cite{zhu2023promptrobust} demonstrates LLMs are fragile under adversarial text perturbations at word and sentence levels.
Pan et al.~\cite{pan2024adversarial} show that semantically aligned adversarial code prompts cause up to 42\% accuracy drops.
Khan et al.~\cite{khan2024cheatingllm} specifically target programming assignment assessments, designing adversarial perturbations that reduce LLM correctness by 77\% to prevent academic dishonesty.
Our approach is complementary: rather than perturbation-based attack, we design a principled \emph{language family} in which LLM priors are systematically misaligned.

\subsection{LLM Impact on Programming Education}
Multiple empirical studies show LLMs reduce independent coding skill development~\cite{dakhel2023github,prather2023robots,lau2023learners}. Novak and McDanel~\cite{novak2025llmhw} explicitly study LLM-resistant assignment design, proposing contextual and personalized tasks; however, they do not characterize \emph{which linguistic features} drive LLM failure.
We address this gap by providing a four-level taxonomy grounded in NTP prior theory.

\subsection{Few-Shot Learning and Pattern Extraction}
In-context learning (ICL) allows LLMs to adapt to new tasks from examples~\cite{dong2022survey}.
Min et al.~\cite{min2022rethinking} show that labels in demonstrations matter less than expected---models rely on surface patterns rather than semantic reasoning.
Our L4 finding (pattern blindness) extends this: LLMs fail to extract \emph{semantic inversion} from worked examples, consistent with the view that ICL is primarily distributional rather than rule-inductive.

\subsection{Research Gap}
Prior adversarial work degrades generation quality without a principled design taxonomy. Prior education work proposes task-level heuristics without mechanism analysis. Prior knowledge-conflict work focuses on factual knowledge rather than programming-language semantics. Our contribution unifies these threads: a theory-grounded taxonomy (L1--L4) that characterizes LLM failure by \emph{type} (prior dominance vs.\ pattern blindness) and \emph{mechanism} (NTP prior conflict), enabling principled LLM-resistant language design.

\subsection{Instruction Alignment}
Alignment methods including RLHF~\cite{ouyang2022instruct} improve instruction following in natural-language tasks, but our results suggest semantic-level programming conflicts remain resistant even in aligned models (gpt-4o succeeds at L3 but fails at L4), pointing to a gap between alignment training and deep semantic inversion under implicit delivery.
